\documentclass[a4paper,12pt]{article}
\usepackage[margin=2cm]{geometry}

% Pacchetti utili di base
\usepackage[utf8]{inputenc}   % codifica dei caratteri
\usepackage[T1]{fontenc}      % font encoding
\usepackage[italian]{babel}   % lingua italiana
\usepackage{graphicx}         % immagini
\usepackage{hyperref}         % link cliccabili
\usepackage{amsmath}          % align center
\usepackage{listings}         %codice programmazione
\usepackage{xcolor}           % per colori opzionali del codice
\usepackage{float}
\usepackage{algorithm}
\usepackage{algpseudocode}
\usepackage{tcolorbox}
\usepackage{tikz}
\usepackage{listings}
\usepackage[utf8]{inputenc}

\usetikzlibrary{arrows.meta, positioning}
\lstset{escapeinside={<@}{@>}}
% Definizione linguaggio MIPS per listings
\lstdefinelanguage{MIPS}{
  morekeywords={lw,sw,add,sub,addi,li,la,j,jr,jal,move,beq,bne,slt},
  sensitive=true,
  morecomment=[l]{\#},      % commento con #
  morestring=[b]"           % stringhe tra doppi apici
}

\newcommand{\datasection}[1]{%
    \section*{#1}%
    \addcontentsline{toc}{section}{#1}%
}

\usepackage{mdframed}

\newenvironment{exambox}{%
    \begin{mdframed}[
        linecolor=black!70,
        linewidth=0.6pt,
        innerleftmargin=8pt,
        innerrightmargin=8pt,
        innertopmargin=6pt,
        innerbottommargin=6pt
    ]
    \begin{description}
}{%
    \end{description}
    \end{mdframed}
}

% Informazioni sul documento
\title{Risposta alle domande G2 di Sistemi Operativi}
\author{Nicolò Giuliani}
\date{}


\begin{document}
\maketitle
\tableofcontents   % <-- genera automaticamente la pagina dell'indice
\newpage 

\datasection{06 febbraio 2026}
\begin{exambox}
  \item[a)] In una tabella FAT possono esistere elementi di uguale valore? (non considerando l’elemento nullo che
contraddistingue la fine del file) Se sì: in quali casi? Se no: quali effetti avrebbe?
  \item[R:] \textit{No, in FAT non possiamo avere elementi di uguale valore, infatti oltre a \texttt{EOF}, che serve per segnare il termine di una catena di puntatori. Se avessimo elementi di ugual valore vorrebbe dire che due puntatori indicano lo stesso blocco. Questo significherebbe che un blocco e' stato sovrascritto da un altro file; ovvero uno o piu' file sono corrotti.}

  \item[b)] Trovare, se possibile, un grafo di holt che abbia 5 nodi, tutte le risorse abbiano molteplicità due (due unità per
ogni classe di risorse) e tutti i processi siano in deadlock
  \item[R:] \textit{Grafo di Holt con 2 processi + 3 classi di risorse (molteplicita' 2) in deadlock.}\\
  \begin{center}
    \begin{tikzpicture}[
  proc/.style={circle, draw, minimum size=1cm, thick},
  res/.style={rectangle, draw, minimum size=1cm, thick},
  >=Stealth, thick
]

\node[proc] (P1) at (0, 1.5)  {$P_1$};
\node[proc] (P2) at (0, -1.5) {$P_2$};

\node[res, label=above:$A$] (A) at (4,  3) {};
\node[res, label=above:$B$] (B) at (4,  0) {};
\node[res, label=above:$C$] (C) at (4, -3) {};

% Dot molteplicità 2
\foreach \x/\y in {4/3, 4/0, 4/-3} {
  \fill (\x-0.2, \y) circle (2.5pt);
  \fill (\x+0.2, \y) circle (2.5pt);
}

% Assegnazioni: risorsa -> processo
\draw[->] (A) -- (P1);
\draw[->] (B) -- (P1);
\draw[->] (C) -- (P1);
\draw[->] (A) -- (P2);
\draw[->] (B) -- (P2);
\draw[->] (C) -- (P2);

% Richieste: processo -> classe
\draw[->, dashed] (P1) to[bend left=15] (B);
\draw[->, dashed] (P2) to[bend right=15] (C);

\end{tikzpicture}
  \end{center}
  \item[c)] Perché i metodi di sincronizzazione tipo test\&set (detti anche spinlock) assumono grande rilevanza nella
scrittura di sistemi operativi multiprocessore?
  \item[R:] \textit{Nei sistemi multiprocessore dobbiamo implementare il kernel in multiprocesso concorrente. Per non avere conflitti tra thread del kernel dobbiamo avere accesso atomico alle strutture dati del Kernel. Per far questo un semplice semaforo non basta dato che blocca a livello solo della singola CPU. Mentre con gli \texttt{spinlock}, ovvero istruzioni dell'IS della CPU, abbiamo accesso atomico su tutti i processori logici contemporaneamente.}
  \label{risposta:spinlock}
  \item[d)] Perché è difficile revocare una autorizzazione fornita tramite capability?
  \item[R:] \textit{Le capability non hanno una "sede centrale" dove possiamo gestirle tutte come per ACL, infatti ogni processo/utente (e anche i figli che la ereditano) ha la sua capability. Quindi in un sistema multiprocesso/multiutente diventa impossibile (o quasi) fare una revoca perche' bisognerebbe scansionare l'intero disco (e sapere anche i figli).}
\end{exambox}
\newpage
\datasection{09 gennaio 2026}
\begin{exambox}
\item[a)] Quali problemi possono accadere se si utilizza un file system ext2 dove un file ha un numero di riferimenti errato?
\item[R:] \textit{Il numero di riferimenti in ext2 indica il numero di hard link del file, un numero errato comporta che abbiamo realmente piu' o meno hardlink.\\
\begin{enumerate}
    \item Caso piu' hardlink: il fs segna che abbiamo piu' hardlink, quindi se andiamo a eliminare l'ultimo \texttt{vero} hardlink il fs non eliminera' internamente l'inode e non lascera' i blocchi come liberi, portando a un inutile memoria occupata per sempre.
    \item Caso meno hardlink: il fs segna che abbiamo meno hardlink, quindi quando andiamo a eliminare un hardlink (anche se realmente ne abbiamo altri) il fs eliminera' l'inode e mettera' i blocchi come liberi, portando a una corruzzione del file.
\end{enumerate} 
Per fixare entrambi i casi fsck (\texttt{File System Consistency Check}) scannerrizza tutti i file e per ogni file conta il numero di hardlink reali e poi corregge se serve nel inode.
} \label{risposta:inode_hardlink}
\item[b)] Siano dati due array di dischi identici, il primo viene gestito in modalità RAID1, il secondo in modalità RAID5. In quale array l'operazione di lettura è statisticamente più veloce?
\item[R:] \textit{Il RAID1 ha come massimo di velocita' in lettura la velocita' di lettura minore di tutti i dischi. Mentre RAID5 avendo avendo oltre che ridondanza anche uno striping abbiamo velocita' in lettura maggiori rispetto anche al singolo disco.}
\item[c)] La seguente descrizione del funzionamento dell'algoritmo del banchiere è errata: ``Quando un processo richiede risorse, l'algoritmo del banchiere controlla se lo stato del sistema rimarrebbe safe assegnando le risorse richieste: se sì le risorse vengono assegnate, se no l'operazione viene sospesa. Quando un processo restituisce risorse si controlla se il nuovo stato consente di soddisfare una o più operazioni sospese.'' Perché è errata? Come si può correggere la descrizione?
\item[R:] \textit{Per la frase \texttt{"... Quando un processo restituisce risorse si controlla se il nuovo stato consente di soddisfare una o più operazioni sospese."} Bisognerebbe dire "... Quando un processo restituisce risorse si riesegue l'algortimo del banchiere per vedere se il nuovo stato e' safe o unsafe." .} 
\item[d)] A cosa serve e come funziona il meccanismo di Aging negli scheduler a priorità dinamica?
\item[R:] \textit{Negli scheduler con priorita' puo' accadere che i processi a bassa priorita' non vengano mai eseguito perche' i processi ad alta priorita' occupano sempre lo scheduler. Quindi per non far accadere questo con il passare del tempo i processi aumentano gradualmente la propria priorita' fino a raggiungere in caso quella massima. Cosi tutti i processi verrano eseguiti.}
\end{exambox}

\newpage
\datasection{05 settembre 2025}
\begin{exambox}
\item[a)] Quando è sconsigliato usare un meccanismo di protezione di accesso di tipo Access Control List?
\item[R:] \textit{Se abbiamo bisogno di semplici permessi \texttt{r/w/x} per \texttt{owner/group/others} i permessi Unix bastano e le ACL aggiungono solo complessita' e overhead. Per i servizi come docker o i file-sharing le ACL non funzionano. E rispetto le capability se vogliamo lavorare su processi e consentire ai processi singoli permessi speciali con capability possiamo mentre con ACL no.}
\item[b)] Qual è la differenza fra una funzione di una libreria ed una system call? In altre parole, perché sono necessarie le system call?
\item[R:] \textit{La libreria generalmente sono fatte in:
\begin{center}
    \boxed{User \, space \rightarrow Syscall \, (kernel \, space) \rightarrow User \, space}
\end{center}
e quindi le librerie sono dei \texttt{wrapper} che ampliano le syscall.\\
Abbiamo pero' sempre bisogno delle syscall infatti queste ci permettono di andare in kernel space per poter accedere alle strutture protette. Questo complessivamente e' un meccanismo di protezione e di agevolamento al programmatore che seguendo la fisolosia della OOP non deve preccuparsi delle implementazioni interne.
}
\item[c)] Il problema del deadlock può essere risolto con tecniche di checkpointing \& rollback?
\item[R:] \textit{Si, pero' questo consiste in snapshot periodoci del sistema, rilevazione runtime del deadlock e ripristino del sistema (recovery). Questo porta come vantaggio il non dover conoscere a priori le risorse di ogni singolo processo rispetto al algortimo del banchiere (preevemption). Pero' porta un overhead di memoria per gli snapshot. Questa tecnica puo' essere ulterioremte ottimizzata nei linguaggi reversibili con modelli:}
\begin{lstlisting}
    try(condition){
        //do something that induce deadlock
    } else{
        //rollback
    }
\end{lstlisting}
\textit{Questo modello non porta overhead di memoria dal rollback dato che non abbiamo snapshot, ma porta a un costo di esecuzione maggiore perche' rilevato un deadlock dobbiamo invertire tutte le operazioni per tornare alla situazione iniziale. Oltre al fatto che tutte le operazioni debbano essere invertibili per natura (a meno che di compilatori che traducono under-the-hood il linguaggio in uno revertibili per poi eseguirlo su interprete).
\begin{align*}
    & f(x) \to f'(x)\\
    & f'(x)^{-1} \to f(x)
\end{align*}
ovvero 
\begin{align*}
    & f(x) \iff f'(x)
\end{align*}
}
\item[d)] Dati $N$ dischi aventi le stesse caratteristiche, il tempo medio di accesso è minore se creiamo con essi un cluster RAID0 o uno RAID5? Discutere la perdita di prestazione per la lettura e la scrittura.
\item[R:] \textit{Con RAID0 abbiamo $N$ accessi, perche' ogni file e' diviso in striping su tutti i dischi. Mentre in RAID5 abbiamo 2 accessi ovvero disco dove scriviamo il blocco e il disco di parita'.\\
\texttt{Lettura}:
\begin{itemize}
    \item RAID0: $N$ accessi paralleli
    \item RAID5: $1$ accesso
\end{itemize}}
\texttt{Scrittura}:
\begin{itemize}
    \item RAID0: $N$ accessi paralleli
    \item RAID5: $4$ accessi sequenziali
    \begin{align*}
        P_{new} = P_{old} \oplus D_{old} \oplus D_{new}
    \end{align*}
    Con $P$ e $D$ rispettivamente per parita' e dato.
\end{itemize}
A livello di tempo medio il RAID0 vince.
\end{exambox}

\newpage
\datasection{21 luglio 2025}
\begin{exambox}
\item[a)] Perché durante la compattazione di memoria occorre sospendere l'esecuzione dei processi coinvolti?
\item[R:] \textit{Perche' i processi portano a frammentazione esterna e quindi si verifica un caso del genere:}\\
Prima:  [P1][---][P2][--][P3]\\
Dopo:   [P1][P2][P3][--------]\\
\textit{Quindi la compattazione sposta i blocchi di memoria del processi in una sequenza contigua per ottimizzazione lo spazio totale. Pero facendo questo gli indirizzi cambiano. Quindi il processo bisogna fermarlo, compattare, calcolare i nuovi indirizzi, applicarlo al processo.}
\item[b)] Perché nell'i-node non è presente il nome del file?
\item[R:] \textit{Nel inode non memorizziamo il nome del file ma il contatore degli hardlink. Infatti se memorizzassimo il nome del file non avremmo la opzione degli hardlink (come per FAT), mentre attarverso il contatore degli hardlink questo e' possibile. Il nome invece e' memorizzato nella directory e anche il respettivo inode.}
\item[c)] Perché il microkernel è meno efficiente del kernel monolitico? Perché è più flessibile e sicuro?
\item[R:] \textit{Il mikrokernel avendo solo il minimo indispensabile in kernelspace e quasi tutti i driver, programmi, ecc in userspace per motivi di flessibilita', passa molto tempo a fare lo swap da usermode a kernelmode (perche' richiamato molte volte) e quindi porta molto overhead. Pero' e' decisamente piu' flessibile perche' se vogliamo modificare un driver non dobbiamo ricompilare tutto il kernel ma solo il singolo programma. E' anche piu' sicuro perche' solo pochi componenti del SO possono accedere alle strutture protette del kernel, rendendo quindi la modifica/lettura molto piu' difficile.}
\item[d)] Nella memorizzazione da parte dei processi utente delle capability per il controllo di accesso può essere usato un algoritmo di crittografia a singola chiave (chiave privata) oppure occorre un algoritmo a doppia chiave (chiave pubblica)? Perché?
\item[R:] \textit{Eslusivamente a chiave privata, perche il kernel firma la capability con la chiave privata (nascosta) e gli utenti/processi controllano la creazione sicura attraverso la loro pubblica. Questo serve perche' se i processi/utenti conoscessero la chiave unica allora potrebbero creare loro una capability e firmarla con essa.}
\end{exambox}

\newpage
\datasection{23 giugno 2025}
\begin{exambox}
\item[a)] L'algoritmo LOOK (detto anche dell'ascensore) è più efficiente di C-LOOK (che ha lo stesso principio di funzionamento del metodo LOOK, ma la scansione del disco avviene in una sola direzione). In quali casi si preferisce C-LOOK a LOOK?
\item[R:]\textit{LOOK ha un thoughtput maggiore ma il tempo tra le chieste agli estremi e' molto alto, mentre C-LOOK va a risolvere il problema andando a minimizzare la varianza (differenza dei tempi medi).}
\item[b)] Spiegare quali attacchi sono possibili se non viene usato il meccanismo del ``sale'' (salt) nella memorizzazione delle password criptate.
\item[R:] \textit{Il \texttt{salt} serve per proteggerci dagli attacchi dizionario. Infatti la password viene confrontata con $H(password + salt)$ con $salt$ come numero univoco per utente. questo numero e' di $256bit \backslash 32byte$.}
\item[c)] L'algoritmo del Banchiere evita il verificarsi di situazioni di deadlock ritardando le richieste che porterebbero il sistema in uno stato unsafe. Basta questa regola a evitare anche situazioni di starvation?
\item[R:] \textit{No, infatti la richiesta potrebbe essere rimandata all'infinito. Bisognerebbe implementare anche una tecnica simile al aging per impedire starvation.2}
\item[d)] Perché è necessario usare spinlock in sistemi multiprocessore per implementare kernel di tipo simmetrico (SMP)?
\item[R:] Vedi risposta a pagina \pageref{risposta:spinlock}, esame 06 febbraio 2026.
\end{exambox}

\newpage
\datasection{28 maggio 2025}
\begin{exambox}
\item[a)] Perché è difficile revocare un diritto di accesso se espresso come capability rispetto a revocare lo stesso diritto in una access control list?
\item[R:] \textit{Perche' in capability non abbiamo una lista centrale per file/dir. come per ACL con tutti i permessi li, ma tante capability per diversi processi/utenti (e anche per i figli). Quindi la revoca diventa impossibile. Mentre per ACL ci basta modificare la tabella riferito a quel file.}
\label{risposta:diff_acl_cap}
\item[b)] Si può implementare la memoria virtuale con paginazione a richiesta in un sistema che ha una Memory Management Unit priva di TLB (Translation Lookaside Buffer)?
\item[R:] \textit{Si, pero' senza la TLB come buffer bisogna ogni volta leggere dalla tabella in RAM. Questo porta' un grande overhead.}
\item[c)] Perché non si può implementare la funzionalità di link fisico di file in un file system di tipo FAT?
\item[R:] \textit{In FAT non esiste un reference counter associato ai blocchi: l'unica informazione presente è la catena nella tabella FAT. Se due directory entry puntassero forzatamente alla stessa catena, al momento della cancellazione di una delle due non ci sarebbe modo di sapere se esistono altri riferimenti, quindi i blocchi verrebbero liberati anche se l'altra entry li sta ancora usando. Per supportare gli hard link servirebbe un contatore di riferimenti, che FAT non prevede.}
\item[d)] L'algoritmo del Banchiere evita che si verifichino situazioni di deadlock. Sembra una funzione desiderabile eppure viene raramente usato nei sistemi reali, perché?
\item[R:] \textit{Perche' nei sistemi reali e' quasi impossibile sapere a priori il numero di risorse che un algortimo ha bisogno, e quindi le condizioni per eseguire l'algortimo diventano impossibili. Poi c'e' anche da tenere in mente il costo computazionale alto perche' per $n$ processi e $r$ valute il costo e' di $O(n^2 \cdot r)$.}
\end{exambox}

\newpage
\datasection{11 febbraio 2025}
\begin{exambox}
\item[a)] La tecnica dell'aging serve per evitare casi di starvation. Come funziona?
\item[R:] \textit{Si, la tecnica del aging serve per gli scheduler a priorita', nei quali se non avessimo aging alcuni processi potrebbero non essere mai eseguiti perche' i processi a priorita' maggiore occupano tutto il tempo la CPU. Quindi con il passare del tempo l'aging aumenta la priorita' di un processo che non e' stato eseguito.}
\item[b)] Può esistere un controller di device capace di usare DMA ma che non genera interrupt? Se NO: perché non esiste? Se SÌ: come deve essere costruito il relativo device driver?
\item[R:] \textit{Puo' esistere ma allora in quel caso dovremmo avere la CPU che fa polling costantemente per vedere quando il trasferimento e' completato. Questo andrebbe ad annullare tutti i benefici del DMA. Quindi gli interrupt ci servono proprio per mettere in pausa la CPU e farla continuare con altri processi che non tocchino gli indirizzi che stiamo lavorando con il DMA.}
\item[c)] Quanti dischi sono coinvolti da una operazione di scrittura di un blocco in un array di dischi RAID 5? Perché?
\item[R:] \textit{In scrittura in RAID5 abbiamo 4 operazioni sequenziali:
\begin{enumerate}
    \item Lettura blocco dati vecchio
    \item Lettura blocco parita' vecchio
    \item Scrittura dati nuovo
    \item Scrittura parita' nuova
\end{enumerate}
\begin{align*}
    P_{new} = P_{old} \oplus D_{old} \oplus D_{new}
\end{align*}
con $P$ e $D$ rispettivamente per parita' e dati.
}
\item[d)] Quale problema risolve la compattazione di memoria? Quali sono i difetti di questo metodo?
\item[R:] \textit{La compattazione di memoria serve per ridurre la frammentazione esterna. Questa puo' essere fatta sia su dischi fissi, ma generalmente su RAM per i processi su CPU senza MMU. Se avessimo una MMU a priori la compatazzione non servirebbe perche' ci basterebbe modificare la tabella degli indirizzi virtuali a indirizzi fisici. Il diffetto principale e' che dobbiamo fermare i processi, spostare i dati, ricalcolare i nuovi indirizzi e copiarli nella PCB del processo e poi riprendere l'esecuzione.}

\end{exambox}

\newpage
\datasection{15 gennaio 2025}
\begin{exambox}
\item[a)] Un i-node di un file system tipo ext2 per un errore viene riportato un numero di link maggiore del reale. Cosa può succedere se si continua ad usare il file system? (perché?) E se il numero di link errato fosse al contrario inferiore al reale cosa potrebbe succedere? (perché?)
\item[R:] Vedi risposta a pagina \pageref{risposta:inode_hardlink}, esame 09 gennaio 2026.
\item[b)] Perché è necessario usare spinlock in sistemi multiprocessore per implementare kernel di tipo simmetrico (SMP)?
\item[R:] Vedi risposta a pagina \pageref{risposta:spinlock}, esame 06 febbraio 2026.
\item[c)] Perché nei sistemi reali l'algoritmo di rimpiazzamento second chance (orologio) viene preferito a LRU sebbene il primo non sia a stack e il secondo sì?
\item[R:] \textit{Nei sistemi reali l'algortimo LRU e' troppo costoso da implementare perche' per avere tempo di accesso alla domanda $O(1$) dovremmo avere una tabella grande $n \times n$, mentre se ci accontentassimo di una risposta $O(n)$ ci basterebbe una vettore $n$ che costantemente anrebbe aggiornato. Il costo di una memoria cosi veloce e di una grandezza cosi esponenziale e' troppo alto in un sistema reale. Mentre per implementare \texttt{second chanche} abbiamo $O(n)$ ma per la memorizzazione solo $n$ bit. E questo algortimo approssima LRU abbastanza precisamente per compensare.}
\item[d)] Perché revocare un'autorizzazione espressa come capability è più difficile che revocare lo stesso diritto quando espresso come access control list?
\item[R:] Vedi risposta a pagina \pageref{risposta:diff_acl_cap}, esame 28 maggio 2025.
\end{exambox}

\newpage
\datasection{10 settembre 2024}
\begin{exambox}
\item[a)] In RAID61 (RAID1 di RAID6) qual è il livello di fault tolerance (quanti dischi possono guastarsi)? E in RAID16 (RAID6 di RAID1)?
\item[R:] {\itshape
In \textbf{RAID61} (RAID1 di RAID6) si mantengono due copie identiche di un array RAID6.
Poiché RAID6 tollera fino a 2 dischi guasti per array, nel caso peggiore garantito
si possono perdere \textbf{almeno 2 dischi qualsiasi}. Nel caso limite, se un intero
array RAID6 cade, l'altro lo copre interamente, quindi si possono perdere tutti i
dischi di un array più 2 dell'altro.

In \textbf{RAID16} (RAID6 di RAID1) i ``dischi'' del RAID6 sono coppie RAID1.
RAID6 tollera la perdita di 2 array interi, e ogni coppia RAID1 tollera 1 disco
guasto internamente. Quindi si possono perdere entrambi i dischi di al massimo
2 coppie (assorbiti dal RAID6), più un disco da ciascuna delle coppie rimanenti
(assorbiti dal RAID1). La fault tolerance scala con il numero di coppie, rendendo
RAID16 generalmente più robusto di RAID61.
}
\item[b)] Nella configurazione di un sistema l'amministratore decide di creare una partizione di un disco rotazionale con file system FAT da montare come \texttt{/tmp}. È una buona o cattiva idea? Perché?
\item[R:] \textit{E' una cattiva idea perche' in FAT per leggere il blocco $n$ dobbiamo prima leggere $n$ blocchi e i file sotto \texttt{/tmp} chiedono molto letture velocemente di blocchi in qualsiasi posizione. Oltretutto per gli HDD peggiorerebbero molto le performance.} 
\item[c)] Quali sono i benefici della gestione degli interrupt nidificati e quali sono le difficoltà da affrontare per sviluppare un sistema operativo che supporti gli interrupt nidificati?
\item[R:] \textit{Il beneficio maggiore e' la maggiore responsivita' del sistema, infatti per gli interrupt di I/O questi possono avere priorita' maggiore rispetto ad altri in corso e questo porta a un sistema piu' responsivo per l'utente. Per far questo ci servirebbe uno stack per poter segnare a quale interrupt tornare, ripristinare il sistema a quando era stato bloccato e non avere cicli infiniti internamente.}
\item[d)] Perché aumentando la dimensione dell'area di memoria secondaria usata dalla memoria virtuale si corre il rischio di thrashing?
\item[R:] \textit{Perche' la CPU accetterebbe l'esecuzione di piu' processi e quindi piu' memoria utilizzata e quindi dovremo passare sempre piu' tempo a passare dati da RAM a memoria secondaria, e avverrebbe trashing peggiorando le performance.}
\end{exambox}
\end{document}